\section{Sincronizacion y fallback manual}

\subsection{Descripcion}

Este documento define la estrategia funcional de sincronizacion automatica de resultados y el uso de fallback manual cuando la integracion externa no sea suficiente. Su objetivo es dejar clara la prioridad de fuentes, las condiciones de override y el disparo controlado de recalculos.

\subsection{Alcance}

Aplica a sincronizaciones manuales, overrides administrativos, reintentos y activacion de recalculo a partir de cambios en resultados o estados de partido.

Nota de alcance para \texttt{QH-111}:
\begin{itemize}
  \item \texttt{QH-111} \textbf{NO} incluye cron programado.
  \item \texttt{QH-111} \textbf{NO} incluye recomputacion real de puntos.
  \item \texttt{QH-111} si incluye caso de uso de sync interno y superficie manual/admin ``Sync now''.
\end{itemize}

\subsection{Definiciones}

\begin{itemize}
  \item \textbf{Sincronizacion automatica}: consulta programada o gatillada a la fuente externa de resultados.
  \item \textbf{Fallback manual}: intervencion administrativa para registrar, corregir o preservar un resultado cuando la fuente automatica no es suficiente.
  \item \textbf{Override manual}: resultado vigente establecido por una accion administrativa autorizada.
\end{itemize}

\subsection{Reglas}

\begin{itemize}
  \item Toda actualizacion de resultado debe ser trazable, idempotente y consistente con el dominio.
  \item Toda correccion manual debe quedar auditada con actor, momento, valor anterior, valor nuevo y justificacion.
  \item Cuando exista un override manual vigente, la sincronizacion automatica no puede sobrescribirlo sin una accion explicita y auditada.
  \item Todo cambio efectivo en resultado vigente o estado relevante del partido debe habilitar un recalculo consistente del scoring y leaderboard; \texttt{QH-111} solo deja el evento listo (no ejecuta recomputacion real).
\end{itemize}

\subsubsection{Secuencia funcional de sincronizacion automatica}

\begin{enumerate}
  \item El sistema consulta al proveedor externo segun su ejecucion manual (QH-111) o ventana operativa definida (ticket posterior).
  \item Cada dato recibido se adapta al modelo interno del dominio.
  \item El sistema determina si el dato es aceptable como candidato para actualizar el resultado vigente.
  \item Si el valor no cambia o no es aplicable, la sincronizacion termina sin alterar scoring.
  \item Si el valor cambia y no existe override manual vigente, el sistema actualiza el resultado vigente, registra evidencia y habilita recalculo.
\end{enumerate}

\subsubsection{Uso del fallback manual}

\begin{itemize}
  \item El fallback manual se usa cuando la fuente externa no entrega el dato a tiempo, lo entrega de forma incompleta, produce inconsistencia o requiere correccion operativa.
  \item La operacion manual debe exigir autorizacion y justificacion visible.
  \item El valor manual pasa a ser la verdad oficial vigente hasta nueva accion explicita y auditada.
  \item El fallback manual no elimina la historia del dato anterior ni la referencia a la fuente de origen.
\end{itemize}

\subsubsection{Eventos que deben habilitar recalculo}

\begin{itemize}
  \item publicacion inicial de resultado oficial utilizable;
  \item correccion de resultado vigente;
  \item cambio de estado de partido que altere scoring o cierre funcional;
  \item override manual que reemplace el valor previamente vigente.
\end{itemize}

\subsubsection{Politica de prioridad de fuentes}

\begin{enumerate}
  \item Si no existe override manual vigente, el sistema puede adoptar el valor aceptado del proveedor externo.
  \item Si existe override manual vigente, ese valor tiene prioridad operativa sobre la API automatica.
  \item Para volver a aceptar automaticamente un valor externo luego de un override manual, debe existir una nueva accion explicita y auditada.
\end{enumerate}

\subsubsection{Ventanas funcionales de sincronizacion}

Para MVP, la estrategia base acordada es corrida batch diaria en ventana operativa fija y disparo manual administrativo cuando se requiera. Funcionalmente el sistema debe cubrir:

\begin{itemize}
  \item sincronizacion diaria por lote de partidos del torneo objetivo;
  \item re-sincronizacion posterior para detectar correcciones de proveedor;
  \item sincronizacion manual administrativa cuando operacion lo requiera.
\end{itemize}

\subsection{Auditoria operativa (QH-111)}

QH-111 debe implementar auditoria operativa del sync. Requerimiento minimo:
\begin{itemize}
  \item \texttt{ExternalResultSyncRun}: corrida (runId, provider, endpoint, parametros, timestamp, outcome).
  \item \texttt{ExternalResultSyncEvent}: evento por match consultado (match externo, match interno si existe, decision aplicada, razon de skip/error, before/after cuando aplique).
\end{itemize}

Notas:
\begin{itemize}
  \item No usar \texttt{AuditLog} como sustituto de estas entidades operativas; \texttt{AuditLog} puede usarse para acciones admin, pero \texttt{ExternalResultSyncRun/Event} son la fuente de verdad operativa del sync.
  \item No crear \texttt{ResultHistory} dentro de QH-111; \texttt{beforeJson/afterJson} en \texttt{ExternalResultSyncEvent} cubre la necesidad del ticket.
\end{itemize}

\subsection{Errores normalizados recomendados}

\begin{itemize}
  \item \texttt{provider\_auth\_invalid}
  \item \texttt{external\_match\_not\_found}
  \item \texttt{external\_match\_not\_mapped}
  \item \texttt{external\_result\_not\_final}
  \item \texttt{manual\_override\_active}
  \item \texttt{provider\_error}
  \item \texttt{provider\_rate\_limited}
  \item \texttt{provider\_payload\_invalid}
\end{itemize}

\subsection{Testing minimo obligatorio (QH-111)}

Documentar como criterio de implementacion que QH-111 cubra al menos:
\begin{itemize}
  \item \texttt{FINISHED} + \texttt{score.fullTime} no null $\rightarrow$ crea \texttt{Result} con \texttt{source=EXTERNAL}.
  \item \texttt{FINISHED} + \texttt{score.fullTime} no null sobre \texttt{Result.source=EXTERNAL} existente $\rightarrow$ actualiza idempotente.
  \item \texttt{TIMED} con score null $\rightarrow$ skip no persistible.
  \item \texttt{LIVE}/\texttt{IN\_PLAY}/\texttt{PAUSED} $\rightarrow$ skip no persistible.
  \item \texttt{POSTPONED}/\texttt{SUSPENDED}/\texttt{CANCELLED} $\rightarrow$ skip no persistible.
  \item \texttt{externalMatchId} no mapeado $\rightarrow$ skip \texttt{external\_match\_not\_mapped}.
  \item \texttt{Result.source=MANUAL\_OVERRIDE} $\rightarrow$ skip \texttt{manual\_override\_active}.
  \item token invalido del proveedor $\rightarrow$ \texttt{provider\_auth\_invalid}.
  \item match inexistente del proveedor $\rightarrow$ \texttt{external\_match\_not\_found}.
  \item batch IDs normaliza correctamente varios matches.
  \item lectura publica refleja resultado sincronizado sin exponer metadata interna.
  \item se crean \texttt{ExternalResultSyncRun} y \texttt{ExternalResultSyncEvent} con evidencia minima.
\end{itemize}

\subsection{Edge Cases}

\begin{itemize}
  \item Si un partido fue postergado antes del lock, la sincronizacion del nuevo kickoff debe reflejarse sin perder consistencia con picks.
  \item Si un partido fue suspendido, la sincronizacion no debe cerrar scoring hasta existir resultado oficial utilizable.
  \item Si una correccion manual reemplaza un dato de API y luego la API vuelve a sincronizar, el valor manual debe preservarse hasta nueva accion explicita.
  \item Si el mismo evento provoca mas de un recalculo, el estado final debe converger al mismo resultado por idempotencia.
\end{itemize}

\subsection{Invariantes}

\begin{itemize}
  \item El fallback manual no elimina trazabilidad.
  \item La integracion automatica no puede sobrescribir un override manual vigente por su sola ejecucion.
  \item Todo recalculo parte del resultado vigente oficial y no de valores intermedios de sincronizacion.
  \item El sistema no debe dejar dos verdades simultaneas activas para un mismo partido.
\end{itemize}

\subsection{Preguntas abiertas}

\begin{itemize}
  \item Confirmar ventanas reales de scheduler, reintentos y re-sincronizacion en ticket posterior de cron (fuera de QH-111).
  \item Confirmar si habra una accion administrativa explicita para liberar un override manual y volver a aceptar sincronizacion automatica normal.
\end{itemize}

\subsection{Decisiones pendientes}

\begin{itemize}
  \item \texttt{Decision pendiente} \texttt{No bloqueante}: decidir si el recalculo posterior a sincronizacion se ejecuta inmediatamente en todos los casos o si algunas rutas usaran cola diferida.
\end{itemize}
